\begin{table*}[t]
\centering
\caption{
    Results of overall evaluation on GrailQA-LeaderBoard (\%).
	% Overall evaluation (EM and F1 on GrailQA-LeaderBoard https://dki-lab.github.io/GrailQA (\%)). 
	\label{tb:overall} 
}
\begin{tabular}{ccccccccc}
\hline
              & \multicolumn{2}{c}{Overall} & \multicolumn{2}{c}{I.I.D.} & \multicolumn{2}{c}{Compositional} & \multicolumn{2}{c}{Zero-Shot} \\ \cline{2-9} 
              & EM           & F1           & EM           & F1          & EM              & F1              & EM            & F1            \\ \hline
%QGG~\cite{lan2020query}           & -            & 36.7         & -            & 40.5        & -               & 33.0            & -             & 36.6          \\ \hline
%S2QL(Anonymous)          & 57.5         & 66.2         & 65.1         & 72.9        & 54.7            & 64.7            & 55.1          & 63.6          \\
%TIARA(Anonymous)         & 73.0         & 78.5         & -            & -           & 68.2            & 76.5            & 68.0         & 73.9         \\ \hline
GrailQA-Rank~\cite{gu2021beyond}  & 50.6         & 58.0         & 59.9         & 67.0        & 45.5            & 53.9            & 48.6          & 55.7          \\
GrailQA-Trans~\cite{gu2021beyond} & 33.3         & 36.8         & 51.8         & 53.9        & 31.0            & 36.0            & 25.7          & 29.3          \\
ReTrack~\cite{chen-etal-2021-retrack}       & 58.1         & 65.3         & 84.4         & 87.5        & 61.5            & 70.9            & 44.6          & 52.5          \\
RNG-KBQA~\cite{ye2022rng}      & 68.8         & 74.4         & 86.2         & 89.0        & 63.8            & 71.2            & 63.0          & 69.2          \\ \hline
\model (Ours)      & \textbf{73.2}            & \textbf{78.7}            & \textbf{88.5}            & \textbf{91.2}           & \textbf{70.0}               & \textbf{76.7}              & \textbf{67.6}             & \textbf{74.0}             \\ \hline
\end{tabular}
\label{tab:overall_eval_GrailQA}
\end{table*}

\subsection{Experimental Settings}
\vpara{Dataset.}
We evaluate our method on GrailQA~\cite{gu2021beyond}, WebQSP~\cite{yih2016value}, and CWQ~\cite{talmor-berant-2018-web}, all of which are based on Freebase. GrailQA focuses on generalization problems which involved up to 4-hop logical expressions and complex operations. WebQSP is an i.i.d. benchmark that required 2-hop reasoning. Although CWQ is not designed to solve generalization problem, we can still separate out the zero-shot test set with all the unseen relations and classes, yielding 576/3519 zero-shot/all test set. 

\vpara{Evaluation Metrics.}
To measure the accuracy of logical expression, we use the well-adopted exact match (EM) which measures the exact equivalence between the query graph of the predicted and the gold logical expression. We also calculate the F1 score based on the predicted and gold answers. 
%To obtain the answers, we build a virtuoso database for Freebase and execute SPARQL queries.

\vpara{Baselines.}
On GrailQA, we mainly compare with the published works on the leaderboard, including GrailQA-Rank~\cite{gu2021beyond}, GrailQA-Trans~\cite{gu2021beyond},  Retrack~\cite{chen-etal-2021-retrack}, RNG-KBQA~\cite{ye2022rng}. They are all SP-based models that target generalization problems in KBQA.
%We do not compare with CBR~\cite{das2021case} and PI Transfer~\cite{cao-etal-2022-program} as their full code is not available. 
On WebQSP and CWQ, we compare our method with the retrieval-based models including GraphNet~\cite{pu2018graphnet},PullNet~\cite{sun2019pullnet} and NSM~\cite{He-WSDM-2021}, and the SP-based models including QGG~\cite{lan2020query}, RNG-KBQA~\cite{ye2022rng}, and PI Transfer~\cite{cao-etal-2022-program}.
We evaluate F1 for the retrieval-based models, while evaluate both F1 and EM for the SP-based methods. We compare all the baselines that have the results on the two datasets or publish the codes that can be executed.


\subsection{Overall Evaluation}
\vpara{Performance.}
In Table~\ref{tb:overall} and Table~\ref{tb:webqsp}, we evaluate the performance of \smodel on different datasets. 
For the baselines, we directly take their results reported in the original papers. To be noted, on the extracted zero-shot test set of CWQ, the results for some models remain empty because their full codes are not released.
As shown in Table~\ref{tab:overall_eval_GrailQA}, our model outperforms all the baselines, especially on the compositional and zero-shot test tasks. Compared with RNG-KBQA, the state-of-the-art published model, we have an absolute gain of 4.3\% and 4.4\% in terms of F1 score and EM respectively. We also outperform on the extracted zero-shot CWQ test set by 11.3\% in terms of F1, as for an unseen complex question, parsing out correct knowledge components and logical skeletons is much easier than directly parsing the coarse-grained logical expression correctly.
% which demonstrates the powerful generalization ability of \smodel on complex questions. 
Since the fine-grained module solely focuses on each component and thus leads to a higher component accuracy, \smodel also outperforms on the i.i.d test set of WebQSP. On the original test set of CWQ, we only under-perform PI Transfer which leverages a pre-train process on a large-scale wiki data that is out scope of CWQ. 

%Although we underperform XXX on all the original test set of CWQ, we can outperform it on the extracted zero-shot test set by xxx\% in terms of F1, which demonstrates the powerful generalization ability of \model. Section~\ref{sec:limitation} examines the failure on some i.i.d. test sets.

%WebQSP has some KB-specific relation-relation pairs which semantic are not close with the question, so the fine-level component detector could fail to cover them. Thus, our model just derives comparable results on WebQSP. We will further improve the performance on i.i.d. dataset in the future work. 

% Since the fine-level component detection is an important 

% step affects the model's learning ability on some KB relevant compositions to some extent, our model does not perform as well as RNG-KBQA on WebQSP. 
% Still, our model give a competitive performance that illustrates the effectiveness.
\begin{table}\centering
\caption{\label{tb:webqsp} F1 Evaluation on WebQSP and CWQ (\%).}
\small
\begin{tabular}{lc|cc}
\hline
 & WebQSP & \multicolumn{2}{c}{CWQ} \\ \cline{2-4} &Overall    & Overall      & Zero-Shot    \\ \hline
GraphNet                                            & 66.4   & 32.8 & 22.3          \\
PullNet & 68.1   & 47.2 & -             \\
NSM    & 74.3   & 48.8 & 31.6          \\ \hline
QGG      & 74.0   & 40.4 & 28.9          \\
RNG-KBQA     & 75.6   & 42.3 & 33.3          \\ 
%KQAPro                                              & 67.5   & 45.9 & 41.8          \\
PI Transfer                        & 76.5   & \textbf{58.7} & -             \\ \hline
Ours                                                & \textbf{76.9}   & 56.4 & \textbf{53.1}          \\ \hline
\end{tabular}
\end{table}			

\vpara{Efficiency.}
%Generally, enumeration process, candidate selection process, and final composition mainly affect the efficiency of each model.
% Generally, the time consuming of enumeration process, candidate selection process, and final composition determine the efficiency of most SP-based KBQA models.
Both RNG-KBQA and GrailQA-Rank enumerate all the logical expressions in a 2-hop KB subgraph (enumeration), so it is time-consuming for the rank model to score thousands of logical expressions for each question (candidate selection). Conversely, our \smodel just retrieves the most relevant components (candidate selection) and then enumerates the component pairs based on the filtered candidates (enumeration), which greatly reduces the inference time. Besides enumeration and candidate selection, a seq-to-seq model is used to generate the final logical expression (final composition).
In the same 24GB GPU and Intel Gold 5218 CPU, the experimental results in Figure~\ref{fig:inference_time} show that our model runs 4 times faster than baselines.

\begin{figure}[t]
\centering
\includegraphics[scale=0.37]{inference_time.png}
\caption{Inference time on GrailQA. ``Overall'' denotes the total inference time of each model. Specially, GrailQA-Rank has no composition step, we record the corresponding time as zero. }     
% \caption{Inference time on GrailQA. The left three pairs of column denotes the inference time of each step. Overall denotes the total inference time of each method. Since GrailQA-Rank does not have composition step, we record the number as zero. }     
\label{fig:inference_time}

\end{figure}


\begin{table*}[t]
\centering
\setlength\tabcolsep{10pt}
\caption{Ablation studies on GrailQA-Dev (\%).}
\label{tb:ablation-result}
\begin{tabular}{lcccccccc}
\hline
& \multicolumn{2}{c}{Overall}   & \multicolumn{2}{c}{I.I.D.}    & \multicolumn{2}{c}{Compositional} & \multicolumn{2}{c}{Zero-Shot} \\ \cline{2-9} 
& EM & F1 & EM  & F1  & EM  & F1  & EM   & F1   \\ \hline
\multicolumn{1}{l}{T5-base}  & 22.7   & 23.4   & 61.8          & 64.1          & 28.3           & 29.0           & 0.3          & 0.3 \\
\multicolumn{1}{l}{RNG-KBQA}  & 71.4   & 76.8   & 86.5          & 88.9          & 61.6            & 68.8            & 69.0          & 74.8          \\ 
\multicolumn{1}{l}{Enhanced RNG-KBQA} & 72.8          & 78.2          & 86.6          & 90.2          & 61.7            & 69.3            & 71.5          & 76.7          \\ \hline
\multicolumn{1}{l}{\model}   & \textbf{79.0} & \textbf{83.8} & \textbf{89.0} & \textbf{91.5} & \textbf{70.4}   & \textbf{77.3}   & \textbf{78.1} & \textbf{83.1} \\ 
\multicolumn{1}{r}{\quad--Knowledge}                   & 23.1          & 24.0          & 62.1          & 64.2          & 29.5            & 31.0            & 0.3          & 0.3          \\   
\multicolumn{1}{r}{\quad--Knowledge Pairs}                  & 53.6          & 55.6          & 70.2          & 72.3          & 44.0            & 46.0            & 50.3          & 52.2   
   \\ 
\multicolumn{1}{r}{\quad--Logical Skeleton}  & 78.0          & 80.8          & 85.2          & 86.8          & 68.5            & 71.9            & 79.2          & 81.8 \\ 
\multicolumn{1}{r}{\quad--Decode Constraint}  & 77.5          & 83.1          & 88.3          & 91.1          & 67.8            & 76.3            & 76.8          & 82.5 \\ \hline
\end{tabular}
\end{table*}

\subsection{Ablation Studies}
\label{sec:ablation}
GrailQA does not provide ground truth for the test set, so we conduct the ablation studies on the public Grail-Dev to investigate how the fine- and middle-grained components affect the performance.

As shown in Table~\ref{tb:ablation-result}, we develop four model variants. (1) \textbf{-Knowledge} removes all the fine-grained and middle-grained components except for the logical skeleton. (2) \textbf{-Knowledge Pairs} replaces the middle-grained components, such as  class-relation pairs and relation-relation pairs with the corresponding fine-grained candidates, such as classes and relations. (3) \textbf{-Logical Skeleton} gets rid of the logical skeleton. 
(4) \textbf{-Decode Constraint} deletes the dynamic vocabulary created with the middle-grained components. 

The results show that  removing ``knowledge'' reduces model performance by 60\% F1 score, and replacing  ``knowledge pairs'' with pure fine-grained components also reduces model performance by 28\% F1, indicating that encoding the middle-grained components can significantly improve the model's reasoning capacity. To further demonstrate that encoding such middle-grained components can also help improve other model's performance, we create Enhanced RNG-KBQA by taking the top-10 ranked results from its ranking model and formulating them into middle-grained component pairs to be injected into its encoder. The results in Table~\ref{tb:ablation-result} show that middle-grained reformulation improves the performance of RNG-KBQA.
Middle-grained component pairs, like coarse-grained logical expressions, can guarantee connectivity, but they are more compact and much shorter.
As a result, because PLMs have a maximum input length, the middle-grained formulation can inject more components and is more likely to cover the components involved in the target logical expression.

Removing ``logical skeleton'' can result in a 3.0\% F1 drop, indicating that skeleton is useful for guiding the question understanding even though it is less important than the knowledge.

Removing ``decode constraint'' in the decoder can also have an effect on model performance, but is much weaker than removing ``knowledge pairs'' in the encoder, indicating that injecting the knowledge constraints in the encoding process is more useful than in the decoding process, because the encoder learns the bidirectional context, which is better suited to natural language understanding. This is also a significant difference from the existing knowledge constrained decoding methods.

Both "Knowledge Pairs" and "Decode Constraint" are proposed for addressing the in-executability issue, which guarantee all generated logical expressions are executable. Removing either reduces the accuracy, which indicates that high executability can improve the model performance.

\subsection{Error Analysis}
We randomly select 50 error cases on GrailQA and summarize the error into three main categories: error entity (60\%), error relation and class (35\%), and error logical skeleton (40\%). We also analysis the error cases while our model fails but some baseline methods can answer successfully resolve them. A typical mistake is on logical expressions that involve KB-specific component composition. For example, in Freebase, ``coach'' is represented by the join of ``sports.sports\_team.coaches'' and ``sports.sports\_team\_coach\_tenure.coach''. Our fine-to-coarse model only predicts the  previous relation but is unable to recall ``sports.sports\_team\_coach\_tenure.coach'', while some coarse-grained methods are able to memorize such composition and provide the correct answer.
